\section*{Hoofdstuk \ref{ch:b1:erosion-transport} --- Erosie en transport}

\begin{solution}{exo:b1:erosion-transport:1}
\qty{0.4}{g.L^{-1}} is \qty{0.4}{kg.m^{-3}}, dus de flux is $500 \times 0.4 =
\qty{200}{kg.s^{-1}}$. Een jaar van \num{365.25}~dagen duurt
\qty{3.156e7}{s}, dus de jaarlijkse flux is \qty{6.31e9}{kg}, of
\qty{6.31}{Mt}.
\end{solution}

\begin{solution}{exo:b1:erosion-transport:2}
$w_s = R g d^2/(18\nu) = 1.6548 \times 9.81 \times (\num{30e-6})^2/(18 \times
\num{1.0034e-6}) = \qty{8.09e-4}{m.s^{-1}}$, ongeveer \qty{0.8}{mm} per seconde. Haar
korrel-Reynoldsgetal is $w_s d/\nu = \num{8.09e-4} \times
\num{30e-6}/\num{1.0034e-6} = 0.024$, ver onder 1: de wet van Stokes geldt. Zij zinkt
door \qty{2}{m} in $2/\num{8.09e-4} = \qty{2472}{s}$, ongeveer
\qty{41}{min}.
\end{solution}

\begin{solution}{exo:b1:erosion-transport:3}
$\tau_b = \rho g h S = 998.2 \times 9.81 \times 0.5 \times \num{4e-4} =
\qty{1.96}{Pa}$. Met $(\rho_s - \rho) g = 1651.8 \times 9.81 =
\qty{16204}{N.m^{-3}}$: voor \qty{1}{mm}, $\theta = 1.96/(16204 \times 0.001) =
0.121$, bijna vier keer de drempel van \num{0.0314}, dus het zand beweegt; voor
\qty{10}{mm}, $\theta = 0.0121$, een vijfde van de drempel van \num{0.0556}, dus het
grind blijft liggen. De stroom voert haar zand over een grindbodem die niet beweegt.
\end{solution}

\begin{solution}{exo:b1:erosion-transport:4}
(a)~$12.6/3.9 = 3.2$: de rivieren voeren ongeveer drie keer meer in suspensie dan in
oplossing. (b)~Grote stuwmeren: $0.20 \times 14 = \qty{2.8}{Gt}$ per jaar; allemaal:
$0.26 \times 14 = \qty{3.64}{Gt}$ per jaar. (c)~De opgeloste flux is
\qty{3.9e15}{g} per jaar bij \qty{120}{g.m^{-3}}, dus het water dat haar meevoert is
$\num{3.9e15}/120 = \qty{3.25e13}{m^3}$, ongeveer \qty{32500}{km^3} per jaar. De
\omterm{def:b1:erosion-transport:load}{zwevende vracht} in dat water is $\qty{1.26e16}{g}/\qty{3.25e16}{L} =
\qty{0.39}{g.L^{-1}}$.
\end{solution}

\begin{solution}{exo:b1:erosion-transport:5}
(a)~De doorsnede heeft het oppervlak $10 \times 2 = \qty{20}{m^2}$ en de natte omtrek
$10 + 2 \times 2 = \qty{14}{m}$, dus $R_h = 20/14 = \qty{1.43}{m}$. (b)~$\tau_b = 998.2 \times 9.81 \times 1.43 \times 0.001 = \qty{14.0}{Pa}$, tegen $\qty{19.6}{Pa}$ voor een zeer
brede geul van \qty{2}{m} diep: de verhouding is $R_h/h = 0.714$, de oevers nemen de rest
van de wrijving voor hun rekening. (c)~Voor een rechthoekige geul met breedte $b$ is
$R_h = bh/(b + 2h) = h/(1 + 2h/b)$, wat naar $h$ neigt wanneer $b$ groeit.
\end{solution}

\begin{solution}{exo:b1:erosion-transport:6}
(a)~$d_1 = \left(18\nu^2/(Rg)\right)^{1/3} = (18/(Rg))^{1/3}\,\nu^{2/3}$, en $R$ en $g$
hangen niet van de viscositeit af. (b)~$\nu = \num{1.3059e-3}/999.7 =
\qty{1.306e-6}{m^2.s^{-1}}$ bij \qty{10}{\celsius}, tegen
\qty{1.0034e-6}{m^2.s^{-1}} bij \qty{20}{\celsius}: de verhouding is \num{1.302}, en
$1.302^{2/3} = 1.192$, dus $d_1 = 0.104 \times 1.192 = \qty{0.124}{mm}$. (c)~In het
gebied van Stokes is $w_s$ omgekeerd evenredig met $\nu$: de korrel bezinkt $1.302$
keer langzamer, ongeveer \qty{30}{\percent}. Koud water houdt fijn sediment langer in
suspensie: een rivier in de winter voert haar silt bij dezelfde afvoer iets verder mee.
\end{solution}

\begin{solution}{exo:b1:erosion-transport:7}
Met $R g/(18\nu) = \qty{8.99e5}{m^{-1}.s^{-1}}$ en de gecombineerde wet:
\begin{center}
\begin{tabular}{@{}lccc@{}}
\hline
$d$ & Stokes (\si{m.s^{-1}}) & gecombineerd (\si{m.s^{-1}}) & verhouding\\
\hline
\qty{0.1}{mm} & \num{8.99e-3} & \num{7.53e-3} & \num{1.19}\\
\qty{0.5}{mm} & \num{0.225} & \num{0.0711} & \num{3.16}\\
\qty{2}{mm} & \num{3.60} & \num{0.197} & \num{18.3}\\
\hline
\end{tabular}
\end{center}
De verhouding groeit met de grootte, omdat het korrel-Reynoldsgetal groeit: de wet van
Stokes, alleen geldig voor kleine Reynoldsgetallen, neemt toe als $d^2$, terwijl de
weerstand van het zog de grove korrels vertraagt tot een snelheid die alleen als $\sqrt d$
groeit.
\end{solution}

\begin{solution}{exo:b1:erosion-transport:8}
(a)~$(Rg/\nu^2)^{1/3} = \qty{25263}{m^{-1}}$. Bij \qty{0.2}{mm} is $D_* = 5.05$
en $\theta_c = 0.3/7.06 + 0.055(1 - e^{-0.101}) = 0.0425 + 0.0053 =
0.0478$; $\tau_c = 0.0478 \times 16204 \times \num{2e-4} = \qty{0.155}{Pa}$.
Bij \qty{20}{mm} is $D_* = 505$ en $\theta_c = 0.3/607 + 0.055(1 - e^{-10.1}) =
0.0005 + 0.0550 = 0.0555$; $\tau_c = 0.0555 \times 16204 \times 0.02 =
\qty{18.0}{Pa}$. (b)~$h = \tau_c/(\rho g S)$ met $\rho g S =
\qty{19.58}{Pa.m^{-1}}$: \qty{7.9}{mm} water zet het fijne zand in beweging, \qty{0.92}{m}
het grind.
\end{solution}

\begin{solution}{exo:b1:erosion-transport:9}
Het bekken verliest \qty{0.8}{Mt} = \qty{8e5}{t} per jaar, of $\num{8e5}/2000 =
\qty{400}{t.km^{-2}.yr^{-1}}$. Volgens \cref{prop:b1:erosion-transport:denudation} is
$D = \qty{8e8}{kg}/(2650 \times \qty{2e9}{m^2}) = \qty{1.51e-4}{m}$ per jaar, of
\qty{151}{m/Myr}. Dat is ongeveer drie keer de mediaan van de met beryllium-10 gemeten
bekkens, \qty{54}{m/Myr}, maar onder hun gemiddelde van \qty{218}{m/Myr}: een tamelijk
snel bekken, zoals in steil of tektonisch actief land.
\end{solution}

\begin{solution}{exo:b1:erosion-transport:10}
(a)~$u_* = \sqrt{g h S} = \sqrt{9.81 \times 3 \times \num{5e-4}} =
\qty{0.1213}{m.s^{-1}}$ en $\tau_b = \rho u_*^2 = \qty{14.7}{Pa}$.
(b)~Suspensie vraagt $w_s < u_* = \qty{0.1213}{m.s^{-1}}$. De gecombineerde wet
geeft \qty{0.126}{m.s^{-1}} bij \qty{1}{mm}; het oplossen van
$w_s(d) = 0.1213$ door proberen geeft $d = \qty{0.94}{mm}$. (c)~Voor grind is
$\theta_c \approx 0.0555$, dus $\tau_c = \qty{14.7}{Pa}$ voor $d \approx
14.7/(0.0555 \times 16204) = \qty{16.3}{mm}$; bij die grootte geeft de drempelkromme
inderdaad $\theta_c = 0.0555$. (d)~Korrels grover dan ongeveer \qty{16}{mm} blijven op de
bodem liggen en vormen een stabiele grindvloer; fijner grind en het zand grover dan
ongeveer een millimeter reizen als \omterm{def:b1:erosion-transport:load}{bodemvracht}; fijner zand, silt en klei reizen in
suspensie.
\end{solution}

\begin{solution}{exo:b1:erosion-transport:11}
(a)~$0.30 \times \qty{1000}{m} = \qty{300}{m}$ gesteente, gemiddeld over de regio.
(b)~De terugveer is $300 \times 2700/3300 = \qty{245}{m}$. (c)~Het gemiddelde oppervlak
is slechts $300 - 245 = \qty{55}{m}$ verlaagd. (d)~De toppen, die niet werden geërodeerd,
worden met de hele regio omhooggedragen: zij staan nu \qty{245}{m} hoger dan voorheen.
Het uitsnijden van dalen kan de toppen van een keten doen stijgen.
\end{solution}

\begin{solution}{exo:b1:erosion-transport:12}
(a)~$0.9 \times \qty{5e8}{kg}/\qty{1300}{kg.m^{-3}} = \qty{346000}{m^3}$ per jaar,
of \qty{0.69}{\percent} van het reservoir. (b)~De helft van het reservoir,
\qty{25e6}{m^3}, is in $\num{25e6}/\num{346000} = \num{72}\text{ jaar}$ gevuld met deze
snelheid. (c)~Naarmate het reservoir zich vult, nemen zijn volume, en dus de verblijftijd
van het water, af, en gaan fijnere korrels erdoor: zijn vangefficiëntie daalt. Ondertussen
worden de oudere afzettingen begraven en gecompacteerd, en wordt water eruit gedreven,
zodat hun volumieke dichtheid stijgt en elke ton minder ruimte inneemt. Beide effecten
vertragen het volumeverlies naarmate het reservoir ouder wordt.
\end{solution}

\begin{solution}{pb:b1:erosion-transport:1}
\textbf{1.} Een jaar water is $1.2 \times \num{3.156e7} = \qty{3.79e7}{m^3}$;
verspreid over \qty{8e7}{m^2}, een laag van \qty{0.473}{m}: een afvoerhoogte van
\qty{473}{mm} per jaar.

\textbf{2.} $1.2 \times \qty{0.150}{kg.m^{-3}} = \qty{0.18}{kg.s^{-1}}$, of
\qty{5680}{t} per jaar.

\textbf{3.} $1.2 \times 0.080 = \qty{0.096}{kg.s^{-1}}$, of \qty{3030}{t} per
jaar.

\textbf{4.} De \omterm{def:b1:erosion-transport:load}{bodemvracht} is \qty{568}{t} per jaar, en de totale $5680 + 3030 +
568 = \qty{9278}{t}$ per jaar.

\textbf{5.} De concentratie van zwevend sediment is hoog gedurende de paar dagen van
vloed, wanneer de afvoer ook hoog is, en laag de rest van de tijd. Een tijdgemiddelde
geeft elke dag hetzelfde gewicht, dus wordt het gedomineerd door de vele dagen van lage
afvoer en lage concentratie. Volgens \cref{prop:b1:erosion-transport:flux} is de
jaarlijkse massa $\int Q\,C\,\mathrm dt$; het door de stroming gewogen gemiddelde is
deze integraal gedeeld door $\int Q\,\mathrm dt$, dus is het door de stroming gewogen
gemiddelde dat met de gemiddelde afvoer moet worden vermenigvuldigd.

\textbf{6.} $3030/9278 = \num{0.33}$, een derde. Voor de rivieren van de wereld,
$3.9/(12.6 + 3.9) = \num{0.24}$, en minder als de \omterm{def:b1:erosion-transport:load}{bodemvracht} werd meegeteld. Het aandeel
van het bekken is hoger: lage heuvels, waar trage mechanische erosie het gesteente de tijd
geeft om op te lossen, en kalksteen en krijt, die volledig oplossen.

\textbf{7.} $9278/80 = \qty{116}{t.km^{-2}.yr^{-1}}$.

\textbf{8.} $D = \qty{9.278e6}{kg}/(2700 \times \qty{8e7}{m^2}) =
\qty{4.30e-5}{m}$ per jaar, \qty{43.0}{m/Myr}.

\textbf{9.} $43.0/54 = 0.80$: iets onder de mediaan van de gemeten bekkens, en vijf keer
onder hun gemiddelde. Een matig langzaam bekken, zoals verwacht voor lage heuvels in een
rustig gebied.

\textbf{10.} De meting averageert over de paar jaren van meting, en hangt af van de
vloeden die daarin plaatsvonden. Beryllium-10 middelt over de tijd die nodig is om de paar
meter te eroderen waarin het wordt geproduceerd: bij \qty{43}{m/Myr} kost elke meter
\num{23000} jaar, dus het gemiddelde beslaat tienduizenden jaren. Zij kunnen verschillen
door zeldzame grote vloeden, die in een korte reeks ontbreken, doordat landbouw en
ontginning de erosie in de afgelopen eeuwen hebben verhoogd, en doordat sediment in de
uiterwaarde wordt opgeslagen en er weer uit wordt gehaald.

\textbf{11.} Een deel van de \omterm{def:b1:erosion-transport:load}{opgeloste vracht} komt niet uit het gesteente: veel van het
bicarbonaat draagt koolstof die uit de lucht of de bodem is opgenomen, en regen brengt
zeezout aan. Ze meetellen zou het bekken krediet geven voor gesteente dat het niet heeft verloren.

\textbf{12.} $\num{1.5e6}/1.2 = \qty{1.25e6}{s}$, \num{14.5} dagen.

\textbf{13.} De gecombineerde wet geeft \qty{7.53e-3}{m.s^{-1}}: de korrel valt
\qty{8}{m} in \qty{1062}{s}, ongeveer \qty{18}{min}. Zij wordt vastgehouden.

\textbf{14.} $w_s = \num{8.99e5} \times (\num{2e-5})^2 =
\qty{3.60e-4}{m.s^{-1}}$: \qty{8}{m} in \qty{22200}{s}, ongeveer \num{6.2}~uur.
Vastgehouden.

\textbf{15.} $w_s = \num{8.99e5} \times (\num{1e-6})^2 =
\qty{8.99e-7}{m.s^{-1}}$: \qty{8}{m} kost \qty{8.9e6}{s}, \num{103}~dagen, zeven keer de
verblijftijd. Het meeste ervan gaat met het water door de dam.

\textbf{16.} De korrel moet \qty{8}{m} in \qty{1.25e6}{s} vallen, dus ten minste
$\qty{6.4e-6}{m.s^{-1}}$. Volgens de wet van Stokes is $d = \sqrt{18\nu
w_s/(Rg)} = \sqrt{\num{6.4e-6}/\num{8.99e5}} = \qty{2.7}{\micro m}$: een fijne klei.
Grovere korrels worden overal waar zij binnenkomen vastgehouden.

\textbf{17.} $568 + 0.7 \times 5680 = \qty{4544}{t}$ per jaar, wat
$\num{4.544e6}/1300 = \qty{3496}{m^3}$ per jaar vult, \qty{0.233}{\percent} van het
reservoir. De helft ervan, \qty{7.5e5}{m^3}, is in ongeveer \num{215} jaar gevuld.

\textbf{18.} $100/43.0 = \qty{2.33}{Myr}$.

\textbf{19.} $1 - 2700/3300 = 0.18$: het oppervlak daalt slechts met
\qty{18}{\percent} van de verwijderde dikte.

\textbf{20.} $100/0.1818 = \qty{550}{m}$ gesteente.

\textbf{21.} De hele regio stijgt met $550 \times 2700/3300 = \qty{450}{m}$ terugveer.
Toppen die nauwelijks zijn geërodeerd, worden ermee omhooggedragen, tot
\qty{450}{m}, terwijl de dalen dieper worden: het reliëf groeit voordat het afneemt.

\textbf{22.} De \omterm{def:b1:erosion-transport:denudation}{afslijtingssnelheid} hangt vooral af van de hellingen. Naarmate de heuvels
worden verlaagd en hun hellingen afvlakken, zal de snelheid dalen, tenzij opheffing of een
klimaatverandering ze weer steiler maakt: de tijd berekend met de snelheid van vandaag is
een ondergrens.

\textbf{23.} $550/42.95 = 12.80$. Bij de huidige snelheid en met isostatische terugveer
kost het verlagen van de Varden Hills met \qty{100}{m} \textbf{12.8}~miljoen jaar.
\end{solution}
